% GENERATED FILE. Edit canonical lessons, then run npm run generate:books.
% canonical-source-hash: fnv1a64:9cff44b1d8f2aacc
% canonical-lessons: FA-C66-hamishe, FA-C66-gahi, FA-C66-zud, FA-C66-dir, FA-C66-ghabl

\chapter{Always, Sometimes, Early, Late, Before}
\label{ch:fa-always-sometimes-early-late-before}
\hlchaptermodality{66}

\begin{chapteropening}
\textbf{By the end of this chapter:} \emph{I can say always, sometimes, early, late and before.}

Complete the last lesson of chapter 66: I can say always, sometimes, early, late and before.
\end{chapteropening}

\section[\texorpdfstring{\fa{همیشه} (hamishe)}{همیشه (hamishe)}]{\fa{همیشه} (hamishe) — always}

\label{lesson:FA-C66-hamishe}



\begin{quote}
Before the new one: say the Persian for tomorrow, then the Persian for now.
\end{quote}

\subsection*{You'll want to know: \fa{همیشه}}
\textbf{\fa{همیشه}} — \emph{hamishe} — \textquotedblleft{}always\textquotedblright{}.

\begin{cousinweb}
\emph{hamishe} is a native Persian word, inherited rather than borrowed; it is usually traced to Middle Persian \emph{hamēšag}, \textquotedblleft{}always\textquotedblright{}.
\end{cousinweb}

\begin{grammarlens}[title={What you've built}]
One more word for this chapter.
\end{grammarlens}

\subsection*{Your turn}
\begin{itemize}
\raggedright
  \item \emph{Say it:} \emph{hamishe}
  \item \emph{Say it:} \emph{hamishe}, once more
  \item \emph{Say it:} say \emph{hamishe}
  \item \emph{Recall:} say the Persian for tomorrow, then the Persian for now, then say \emph{hamishe} again
\end{itemize}

\begin{tcolorbox}[breakable,colback=teal!4,colframe=teal!35!black,title={Before you move on}]
What does \fa{همیشه} mean? (\textquotedblleft{}Always\textquotedblright{}.) Say it once more.
\end{tcolorbox}
\section[\texorpdfstring{\fa{گاهی} (gâhi)}{گاهی (gâhi)}]{\fa{گاهی} (gâhi) — sometimes}

\label{lesson:FA-C66-gahi}



\begin{quote}
Before the new one: say the Persian for now, then the Persian for always.
\end{quote}

\subsection*{You'll want to know: \fa{گاهی}}
\textbf{\fa{گاهی}} — \emph{gâhi} — \textquotedblleft{}sometimes\textquotedblright{}.

\begin{cousinweb}
\emph{gâhi} is \emph{gâh}, \textquotedblleft{}time, place\textquotedblright{}, with the ending \emph{-i}, \textquotedblleft{}a\textquotedblright{}: \textquotedblleft{}a time\textquotedblright{}, so \textquotedblleft{}sometimes\textquotedblright{}. \emph{gâh} is native and usually traced to Middle Persian \emph{gāh}.
\end{cousinweb}

\begin{grammarlens}[title={What you've built}]
One more word for this chapter.
\end{grammarlens}

\subsection*{Your turn}
\begin{itemize}
\raggedright
  \item \emph{Say it:} \emph{gâhi}
  \item \emph{Say it:} \emph{gâhi}, once more
  \item \emph{Say it:} say \emph{gâhi}
  \item \emph{Recall:} say the Persian for now, then the Persian for always, then say \emph{gâhi} again
\end{itemize}

\begin{tcolorbox}[breakable,colback=teal!4,colframe=teal!35!black,title={Before you move on}]
What does \fa{گاهی} mean? (\textquotedblleft{}Sometimes\textquotedblright{}.) Say it once more.
\end{tcolorbox}
\section[\texorpdfstring{\fa{زود} (zud)}{زود (zud)}]{\fa{زود} (zud) — early}

\label{lesson:FA-C66-zud}



\begin{quote}
Before the new one: say the Persian for always, then the Persian for sometimes.
\end{quote}

\subsection*{You'll want to know: \fa{زود}}
\textbf{\fa{زود}} — \emph{zud} — \textquotedblleft{}early\textquotedblright{}.

\begin{cousinweb}
\emph{zud} is a native Persian word, inherited rather than borrowed; it is usually traced to Middle Persian \emph{zūd}, \textquotedblleft{}quick\textquotedblright{}.
\end{cousinweb}

\begin{grammarlens}[title={What you've built}]
One more word for this chapter.
\end{grammarlens}

\subsection*{Your turn}
\begin{itemize}
\raggedright
  \item \emph{Say it:} \emph{zud}
  \item \emph{Say it:} \emph{zud}, once more
  \item \emph{Say it:} say \emph{zud}
  \item \emph{Recall:} say the Persian for always, then the Persian for sometimes, then say \emph{zud} again
\end{itemize}

\begin{tcolorbox}[breakable,colback=teal!4,colframe=teal!35!black,title={Before you move on}]
What does \fa{زود} mean? (\textquotedblleft{}Early\textquotedblright{}.) Say it once more.
\end{tcolorbox}
\section[\texorpdfstring{\fa{دیر} (dir)}{دیر (dir)}]{\fa{دیر} (dir) — late}

\label{lesson:FA-C66-dir}



\begin{quote}
Before the new one: say the Persian for sometimes, then the Persian for early.
\end{quote}

\subsection*{You'll want to know: \fa{دیر}}
\textbf{\fa{دیر}} — \emph{dir} — \textquotedblleft{}late\textquotedblright{}.

\begin{cousinweb}
\emph{dir} is a native Persian word, inherited rather than borrowed; it is usually traced to Middle Persian \emph{dēr}, \textquotedblleft{}long, late\textquotedblright{}. Old Persian \emph{darga-}, \textquotedblleft{}long\textquotedblright{}, is usually given as its older form. (EtymDB also lists a word for \textquotedblleft{}monastery\textquotedblright{} spelled the same, which is unrelated.)
\end{cousinweb}

\begin{grammarlens}[title={What you've built}]
One more word for this chapter.
\end{grammarlens}

\subsection*{Your turn}
\begin{itemize}
\raggedright
  \item \emph{Say it:} \emph{dir}
  \item \emph{Say it:} \emph{dir}, once more
  \item \emph{Say it:} say \emph{dir}
  \item \emph{Recall:} say the Persian for sometimes, then the Persian for early, then say \emph{dir} again
\end{itemize}

\begin{tcolorbox}[breakable,colback=teal!4,colframe=teal!35!black,title={Before you move on}]
What does \fa{دیر} mean? (\textquotedblleft{}Late\textquotedblright{}.) Say it once more.
\end{tcolorbox}
\section[\texorpdfstring{\fa{قبل} (ghabl)}{قبل (ghabl)}]{\fa{قبل} (ghabl) — before}

\label{lesson:FA-C66-ghabl}



\begin{quote}
Before the new one: say the Persian for early, then the Persian for late.
\end{quote}

\subsection*{You'll want to know: \fa{قبل}}
\textbf{\fa{قبل}} — \emph{ghabl} — \textquotedblleft{}before\textquotedblright{}.

\begin{cousinweb}
\emph{ghabl} is borrowed from Arabic \emph{qabl}, \textquotedblleft{}before\textquotedblright{}. Its Arabic root is \emph{q-b-l} (Lane).
\end{cousinweb}

\begin{grammarlens}[title={What you've built}]
One more word for this chapter.
\end{grammarlens}

\subsection*{Your turn}
\begin{itemize}
\raggedright
  \item \emph{Say it:} \emph{ghabl}
  \item \emph{Say it:} \emph{ghabl}, once more
  \item \emph{Say it:} say \emph{ghabl}
  \item \emph{Recall:} say the Persian for early, then the Persian for late, then say \emph{ghabl} again
\end{itemize}

\begin{tcolorbox}[breakable,colback=teal!4,colframe=teal!35!black,title={Before you move on}]
What does \fa{قبل} mean? (\textquotedblleft{}Before\textquotedblright{}.) Say it once more.
\end{tcolorbox}
